% Simeon Nedelchev — CV (Russian)
% Compile: cd cv && latexmk -pdf cv_ru.tex && cp cv_ru.pdf ../pdf/CV_ru.pdf

\documentclass[11pt,a4paper]{article}

\usepackage[margin=0.7in]{geometry}
\usepackage[T1,T2A]{fontenc}
\usepackage[utf8]{inputenc}
\usepackage[russian]{babel}
\usepackage{cmap}
\usepackage{microtype}
\usepackage{enumitem}
\usepackage{titlesec}
\usepackage[hidelinks]{hyperref}
\usepackage{tabularx}

\pagestyle{empty}
\setlength{\parindent}{0pt}
\setlength{\parskip}{0.35em}

\titleformat{\section}{\large\bfseries}{}{0pt}{}[\titlerule]
\titlespacing*{\section}{0pt}{0.85em}{0.4em}

\newcommand{\cvheader}[1]{{\bfseries #1}}
\newcommand{\cventry}[4]{%
  \textbf{#1} \hfill \textit{#2}\\
  \textit{#3}\\
  #4\par\vspace{0.25em}
}
\newcommand{\cvitem}[2]{\item #1 #2}

\hypersetup{
  pdftitle={Simeon Nedelchev — CV (Russian)},
  pdfauthor={Simeon Nedelchev},
}

\begin{document}

\begin{center}
  {\LARGE\bfseries SIMEON NEDELCHEV}\\[0.35em]
  \href{mailto:simkaned@gmail.com}{simkaned@gmail.com} \enspace$\cdot$\enspace
  \href{mailto:s.nedelchev@innopolis.university}{s.nedelchev@innopolis.university}\\[0.15em]
  \href{https://simeon-ned.github.io}{simeon-ned.github.io} \enspace$\cdot$\enspace
  \href{https://simeon-ned.github.io/presentations/template/}{Презентация} \enspace$\cdot$\enspace
  \href{https://github.com/simeon-ned}{github.com/simeon-ned} \enspace$\cdot$\enspace
  \href{https://scholar.google.com/citations?user=Jl8IxhgAAAAJ}{Google Scholar}
\end{center}

\vspace{0.4em}

\section{Профессиональный опыт}

\cventry{Исследователь и старший RL-инженер}{февр. 2026 -- н.в.}{Институт искусственного интеллекта}{%
  Разработка обучаемого управления всем телом, генерации движений на основе данных и методов оценки состояния для гуманоидных и артикулированных робототехнических систем. Построение симуляционных конвейеров, библиотек движений и развёртываемых стеков управления в сотрудничестве с академическими и отраслевыми партнёрами.}

\cventry{Старший инженер по оптимальному управлению и RL}{авг. 2023 -- дек. 2025}{Центр робототехники Сбер, Сбер}{%
  Разработчик RL-политик управления всем телом (WBC) и локомоции, продемонстрированных на гуманоиде Green на конференции AIJ 2025. Работа сфокусирована на управлении, моделировании и симуляции в MuJoCo, конвейерах обучения и оценки; частичная идентификация системы (в основном параметры двигателей и приводов).%
  (\href{https://youtu.be/eehll8Uw3pw?t=1432}{видео}).}

\cventry{Старший преподаватель}{апр. 2022 -- н.в.}{Институт робототехники и компьютерного зрения, Университет Иннополис}{%
  Продвинутые курсы по робототехнике, математическому моделированию, планированию движения, прикладному и нелинейному управлению. Разработка учебных программ и научное руководство студентами в области динамики, симуляции и embodied AI.}

\cventry{Младший научный сотрудник}{сент. 2019 -- сент. 2023}{Центр технологий в робототехнике и мехатронных компонентах, Университет Иннополис}{%
  Математическое моделирование, идентификация и управление приводами на скрученных тросах; оптимальное нелинейное управление и оптимизация траекторий; механическое и электрическое прототипирование новых робототехнических систем.}

\cventry{Научный сотрудник}{мар. 2017 -- июн. 2019}{Лаборатория биоробототехники, Корейский технологический университет образования}{%
  Анализ, прототипирование и нелинейное управление приводами на скрученных тросах (TSA); аппаратная реализация алгоритмов продвинутого управления.}

\section{Образование}

\cventry{Аспирантура по компьютерным наукам (робототехника), обучение завершено}{июл. 2023}{Университет Иннополис}{Средний балл 4{,}87/5 (97/100). Институт робототехники и компьютерного зрения.\\
Диссертация: \emph{Dynamic parameter estimation and adaptive control over robotic systems: generalized momentum approach}.}

\cventry{M.Sc.\ по машиностроению}{февр. 2019}{Корейский технологический университет образования}{Средний балл 4{,}45/4{,}5 (98/100).\\
Диссертация: \emph{Design of Robotic Gripper with Constant Transmission Ratio Based on Twisted String Actuator}.}

\cventry{M.Sc.\ по робототехнике и мехатронике}{июн. 2018}{Московский государственный технологический университет «СТАНКИН»}{Средний балл 4{,}9/5 (98/100).\\
Диссертация: \emph{The Dynamics of Controlled Motion of Industrial Robots}.}

\cventry{B.Sc.\ по робототехнике}{июн. 2016}{Московский государственный технологический университет «СТАНКИН»}{Средний балл 4{,}35/5 (87/100).\\
Диссертация: \emph{Development and research of dynamic model of robot manipulator}.}

\section{Избранные исследовательские и программные проекты}

\textbf{WBC-Gen: Reactive Generative Motion Planner}\\
Быстрый проприоцептивно обусловленный latent flow-matching планировщик, выдающий отслеживаемые WBC-референсы для замороженного контроллера всего тела; обобщается на разные задачи на одном WBC-контракте; способен \href{https://youtu.be/gAHCAtJnRDU}{автономно восстанавливаться и вставать из failed states} без FSM режимов fall/standing. Gen и WBC совместно на onboard CPU Unitree G1, с \href{https://youtube.com/shorts/dXNI8benYAo}{walk/sprint телеопом}, \href{https://youtube.com/shorts/iofTKa0GPwQ}{отклонением возмущений} и \href{https://youtube.com/shorts/aMezmX262tg}{agile recovery} на железе. \href{https://wbc-mjlab.github.io/wbc-demo/}{Попробуйте Gen вместе с WBC-трекингом}, чтобы интерактивно генерировать локомоцию робота (WASD — управление, Shift — спринт; G переключает clips~$\leftrightarrow$~Gen); код деплоя в \href{https://github.com/wbc-mjlab/wbc-g1-deploy/tree/feat/generator}{wbc-g1-deploy}; training code и статья скоро.

\textbf{WBC-Mjlab: Whole Body Control in MuJoCo Lab} (\href{https://github.com/wbc-mjlab/wbc-mjlab}{github.com/wbc-mjlab/wbc-mjlab}) — \href{https://wbc-mjlab.github.io/wbc-demo/}{веб-демо}\\
Общий MDP для универсального отслеживания движений всего тела в mjlab: обучение на нескольких клипах, пресеты задач из статей (ZEST, BeyondMimic-style RSI и др.), конвейер данных движений и экспорт sim-to-real для гуманоидов.

\textbf{PSM: Predictive Style Matching for Natural Locomotion} (\href{https://simeon-ned.github.io/predictive-style-matching/}{страница проекта})\\
Офлайн-предиктор сопоставляет состояние нижней части тела и команды скорости с целями верхней части тела и походки при RL-обучении; развёрнут на Unitree G1 с инференсом только по задаче в runtime. arXiv:2606.07083 (2026).

\textbf{Locomotion and Whole Body Control for Sber Green} (\href{https://youtu.be/eehll8Uw3pw?t=1432}{презентация AIJ 2025})\\
Разработчик RL-политик локомоции и управления всем телом, продемонстрированных на гуманоиде Green на AIJ 2025: повышение естественности и «человекоподобности» ходьбы; RL-based WBC, обученный на синтетически сгенерированных опорных траекториях (без mocap-датасетов и ретаргетинга), и единое отслеживание для телеуправления, танцев и zero-shot навыков.

\section{Вклад в open-source проекты}

\textbf{Pinocchio} (\href{https://github.com/stack-of-tasks/pinocchio}{stack-of-tasks/pinocchio}) — участие в библиотеке динамики твёрдых тел (параметризации инерции и физическая согласованность).

\textbf{Pink} (\href{https://github.com/stephane-caron/pink}{stephane-caron/pink}) — участие в дифференциальной обратной кинематике для Pinocchio (задачи, ограничения, примеры).

\textbf{Mink} (\href{https://github.com/kevinzakka/mink}{kevinzakka/mink}) — вклад \texttt{EqualityConstraintTask} для замкнутых кинематических цепей; примеры для гуманоидов и манипуляции в MuJoCo-порте Pink.

\textbf{GMR: General Motion Retargeting} (\href{https://github.com/YanjieZe/GMR}{YanjieZe/GMR}) — участие (формат SOMA BVH и конфигурации для нескольких роботов).

\textbf{mujoco-sysid} (\href{https://github.com/based-robotics/mujoco-sysid}{based-robotics/mujoco-sysid}) — участие; идентификация систем в MuJoCo и MJX (оценка параметров, регрессоры динамики).

\textbf{MJINX: Differentiable GPU-accelerated Numerical IK} (\href{https://github.com/based-robotics/mjinx}{based-robotics/mjinx}) — соавтор (JAX и MuJoCo MJX).

\section{Публикации}

\textbf{Препринты}
\begin{itemize}[leftmargin=*, nosep, topsep=0.25em]
  \cvitem{}{\textbf{Nedelchev S}, Chaikovskaia E, Davydenko E, Zaliaev E, Gorbachev R. Predictive Style Matching: Natural and Robust Humanoid Locomotion. arXiv:2606.07083, 2026.}
  \cvitem{}{Maslennikov E, Zaliaev E, Dudorov N, Shamanin O, Karanov D, Afanasev G, Burkov A, Lygin E, \textbf{Nedelchev S}, Ponomarev E. Robust RL Control for Bipedal Locomotion with Closed Kinematic Chains. arXiv:2507.10164, 2025.}
  \cvitem{}{Alentev I, Kozlov L, Domrachev I, \textbf{Nedelchev S}, Ryu JH. VIMPPI: Enhancing Model Predictive Path Integral Control with Variational Integration for Underactuated Systems. arXiv:2505.05507, 2025.}
\end{itemize}

\textbf{Статьи в журналах}
\begin{itemize}[leftmargin=*, nosep, topsep=0.25em]
  \cvitem{}{\textbf{Nedelchev S}, Kozlov L, Khusainov RR, Gaponov I. Enhanced Adaptive Control over Robotic Systems via Generalized Momentum Dynamic Extensions. \emph{Russian Journal of Nonlinear Dynamics}, 19(4):633--646, 2023.}
  \cvitem{}{Fam CA, \textbf{Nedelchev S}. Optimization Driven Robust Control of Mechanical Systems with Parametric Uncertainties. \emph{Russian Journal of Nonlinear Dynamics}, 19(4):585--597, 2023.}
  \cvitem{}{Skvortsova V, \textbf{Nedelchev S}, Brown J, Farkhatdinov I, Gaponov I. Design, characterisation and validation of a haptic interface based on twisted string actuation. \emph{Frontiers in Robotics and AI}, 2022.}
  \cvitem{}{\textbf{Nedelchev S}, Skvortsova V, Guryev B, Gaponov I, Ryu JH. On Energy-Preserving Motion in Twisted String Actuators. \emph{IEEE Robotics and Automation Letters}, 2021.}
  \cvitem{}{\textbf{Nedelchev S}, Gaponov I, Ryu JH. Accurate Dynamic Modeling of Twisted String Actuators Accounting for String Compliance and Friction. \emph{IEEE RA-L}, 2020.}
\end{itemize}

\textbf{Конференционные доклады}
\begin{itemize}[leftmargin=*, nosep, topsep=0.25em]
  \cvitem{}{Jnadi A, Khusainov RR, \textbf{Nedelchev S}, Savin S. Explicit Model Predictive Control Design based on Constrained Zonotope Propagation. \emph{DCNA}, 2023.}
  \cvitem{}{\textbf{Nedelchev S}, Kirsanov D, Gaponov I, Seong H, Ryu JH. On Smooth Time-Optimal Trajectory Planning in Twisted String Actuators. \emph{ICRA}, 2021.}
  \cvitem{}{Sabirova A, \textbf{Nedelchev S}, Gaponov I. Parameter Identification in Mechanical Systems with Energy-Based Regressor: Preliminary Study. \emph{NIR}, 2021.}
  \cvitem{}{\textbf{Nedelchev S}, Kirsanov D, Gaponov I. IMU-based Parameter Identification and Position Estimation in Twisted String Actuators. \emph{IROS}, 2020.}
  \cvitem{}{Balakhnov O, \textbf{Nedelchev S}, Gaponov I. Preliminary Study on Slack-Free MPC of Twisted String-Based Antagonistic Joints. \emph{NIR}, 2020.}
  \cvitem{}{\textbf{Nedelchev S}, Gaponov I, Ryu JH. High-Bandwidth Control of Twisted String Actuators. \emph{ICRA}, 2019.}
  \cvitem{}{\textbf{Nedelchev S}, Gaponov I, Ryu JH. Design of Robotic Gripper with Constant Transmission Ratio Based on Twisted String Actuator. \emph{IROS}, 2018.}
  \cvitem{}{Kosterev D, Vorotnikov A, \textbf{Nedelchev S}, Romash E, Poduraev Y. Development of 2-DoF Adaptive Mechatronic Device with Corrective Adjustment of Laser Tracker Reflector for Industrial Robot Calibration. \emph{Annals of DAAAM \& Proceedings}, 28, 2017.}
\end{itemize}

\section{Преподавание}

\textbf{Университет Иннополис} — старший преподаватель (2022--н.в.), ассистент преподавателя (2019--2022)

\textbf{FORC: Fundamentals of Robot Control} (\href{https://github.com/simeon-ned/forc}{github.com/simeon-ned/forc})\\
Вводный курс по моделированию в пространстве состояний, устойчивости, линейному и нелинейному управлению, линеаризации обратной связью для робототехнических систем (лекции и лабораторные в Colab/Python).

\textbf{MCP: Modern Control Paradigms} (\href{https://github.com/simeon-ned/mcp}{github.com/simeon-ned/mcp})\\
Материалы курса по современному управлению через численные методы и выпуклую оптимизацию с интерактивными ноутбуками.

\begin{itemize}[leftmargin=*, nosep, topsep=0.2em]
  \item Прикладное нелинейное управление; современные парадигмы управления; основы управления роботами
  \item Моделирование и симуляция робототехнических систем; теория линейного управления; робототехнические системы
  \item Вычислительный интеллект; продвинутая робототехника
\end{itemize}

\section{Технические навыки}

\textbf{Программирование:} Python, C/C++, Bash\\
\textbf{Робототехника и ML:} MuJoCo, mjlab, Isaac Sim, Isaac Lab, конвейеры обучения RL/WBC, идентификация систем\\
\textbf{Встроенные системы:} STM32, ARM, ESP32, RP2040 (MicroPython, HAL, FreeRTOS, mbed)\\
\textbf{Инструменты:} Linux, LaTeX, Docker, uv, pixi

\section{Награды и достижения}

\begin{itemize}[leftmargin=*, nosep, topsep=0.2em]
  \item Лучшая магистерская диссертация, Корейский технологический университет образования, 2019
  \item Первая премия, всероссийский конкурс дипломных работ «Be-First» (компьютерные и информационные технологии), 2018
\end{itemize}

\section{Научные интересы}

Сочленённые роботы; обучение с подкреплением; гибридные методы; sim-to-real; имитация движений; управление всем телом и локомоция гуманоидов; обучаемое и нелинейное управление; идентификация и оценка состояния систем; дифференцируемая симуляция; генерация движений на основе данных.

\end{document}
